% Source: https://artofproblemsolving.com/community/c3h1279588_graph_of_functions
% Stable identity: aops-topic-1279588-q01
Reflecting $y=f(x)$ in the $y$-axis gives $y=f(-x)$. Shifting a graph $y=g(x)$ two units left gives $y=g(x+2)$, so from $g(x)=f(-x)$ we get $y=f(-(x+2))=f(-x-2)$, which is not $f(2-x)$ in general.

Correct descriptions, since $f(2-x)=f(-(x-2))$:
\begin{itemize}
\item reflect in the $y$-axis to get $f(-x)$, then shift $2$ units to the \emph{right} to get $f(-(x-2))=f(2-x)$; or
\item shift $2$ units left to get $f(x+2)$, then reflect in the $y$-axis to get $f(-x+2)$.
\end{itemize}
Both amount to a single reflection in the vertical line $x=1$: the point $(a,b)$ on $y=f(x)$ goes to $(2-a,b)$, and indeed $f(2-(2-a))=f(a)=b$.
